%% kapa: exam=2 round=29 year=2018 period=1 subject=appraisal_practice subject_ko=감정평가실무 seq=2 number=2 points=30
%% source: Q-Net 감정평가사 기출문제 (cst00309) artlSeq=5203398 file=raw/1008/5203398/1교시 감정평가실무.pdf
%% conversion: kordoc 4.24.51 -> tools/gen2.py (auto); review=proofread-sonnet; flags=none
\begin{longquestion}{【문제 2】}
감정평가사 甲은 ○○공사로부터 소유건축물의 일부(1층 101호)에 대해 2018. 07. 01. 기준시점의 임대료감정평가를 의뢰받았다. 관련 법규와 감정평가이론을 참작하고, 제시된 자료를 활용하여 다음의 물음에 답하시오. (30점) \par
\medskip\noindent 물음 1) 적산법에 의한 임대료를 산정하시오. (20점) \par
\medskip\noindent 물음 2) 임대사례비교법에 의한 임대료를 산정하고, 적산법에 의한 임대료와 시산 가액조정을 통해 임대료를 결정하시오. (10점) \par
< 자료 1 > 대상물건의 토지 내역 \par
1. 소재지: A광역시 S구 S동 118번지 \par
2. 용도지역: 근린상업지역 \par
3. 대상토지 현황: 대, 350 m\textsuperscript{2}, 광대한면, 부정형, 평지 \par
4. 주위환경: 대로변을 따라 5층 ∼ 10층 규모의 금융회사, 사무실 등 상업용 또는 업무용건물이 밀집하여 위치함. 후면은 소로를 따라 저층규모의 주상용건물, 일부 단독주택 등이 혼재하고 있음 \par
< 자료 2 > 대상물건의 건물 내역 \par
1. 건물 현황: 건축면적 250 m\textsuperscript{2}, 연면적 1,740 m\textsuperscript{2} \par
2. 건물구조 등: 철근콘크리트조 슬래브지붕 지하 1층, 지상 6층 \par
3. 사용승인일: 2000. 05. 08. \par
4. 층별 현황(집합건축물대장) \par
(단위: m\textsuperscript{2}) \par
{\small\centering
\begin{tabular}{|c|c|c|c|c|c|}
\hline
\makecell{층별} & \makecell{용도} & \makecell{바닥면적} & \makecell{전유면적} & \makecell{공유면적} & \makecell{비고} \\
\hline
\makecell{지하1층} & \makecell{업무시설} & \makecell{300} & \makecell{220} & \makecell{80} &  \\
\hline
\makecell{1층} & \makecell{업무시설, 소매점} & \makecell{250} & \makecell{188} & \makecell{62} & \makecell{3개 호실} \\
\hline
\makecell{2층} & \makecell{업무시설} & \makecell{250} & \makecell{188} & \makecell{62} &  \\
\hline
\makecell{3층} & \makecell{업무시설} & \makecell{240} & \makecell{180} & \makecell{60} &  \\
\hline
\makecell{4층} & \makecell{업무시설} & \makecell{240} & \makecell{180} & \makecell{60} &  \\
\hline
\makecell{5층} & \makecell{업무시설} & \makecell{240} & \makecell{180} & \makecell{60} &  \\
\hline
\makecell{6층} & \makecell{업무시설} & \makecell{220} & \makecell{165} & \makecell{55} &  \\
\hline
\end{tabular}\par}

5. 1층 호별 현황(집합건축물대장) \par
(단위: m\textsuperscript{2}) \par
{\small\centering
\begin{tabular}{|c|c|c|c|}
\hline
\makecell{1층} & \makecell{101호} & \makecell{102호} & \makecell{103호} \\
\hline
\makecell{전유면적} & \makecell{60} & \makecell{55} & \makecell{73} \\
\hline
\end{tabular}\par}

< 자료 3 > 표준지공시지가 현황(공시기준일: 2018. 01. 01.) \par
{\small\centering\footnotesize\setlength{\tabcolsep}{3pt}
\begin{tabular}{|c|c|c|c|c|c|c|c|c|c|}
\hline
\makecell{기호} & \makecell{소재지} & \makecell{지목} & \makecell{면적\newline{}(m\textsuperscript{2})} & \makecell{이용상황} & \makecell{용도지역} & \makecell{도로\newline{}교통} & \makecell{형상\newline{}지세} & \makecell{공시지가\newline{}(원/m\textsuperscript{2})} & \makecell{비고} \\
\hline
\makecell{1} & \makecell{S구 S동\newline{}9} & \makecell{대} & \makecell{588.0} & \makecell{상업용} & \makecell{근린상업} & \makecell{광대한면} & \makecell{세장형\newline{}평지} & \makecell{4,300,000} & \makecell{계획도로\newline{}20 \% 저촉} \\
\hline
\makecell{2} & \makecell{S구 S동\newline{}40} & \makecell{대} & \makecell{388.0} & \makecell{주상용} & \makecell{근린상업} & \makecell{중로\newline{}한면} & \makecell{사다리\newline{}평지} & \makecell{3,600,000} &  \\
\hline
\makecell{3} & \makecell{S구 S동\newline{}261} & \makecell{대} & \makecell{550.0} & \makecell{업무용} & \makecell{3종\newline{}일주} & \makecell{광대\newline{}소각} & \makecell{가장형\newline{}평지} & \makecell{4,000,000} &  \\
\hline
\end{tabular}\par}

< 자료 4 > 평가사례 \par
{\small\centering\footnotesize\setlength{\tabcolsep}{3pt}
\begin{tabular}{|c|c|c|c|c|c|c|c|c|c|c|}
\hline
\makecell{기호} & \makecell{소재지} & \makecell{지목} & \makecell{면적\newline{}(m\textsuperscript{2})} & \makecell{이용\newline{}상황} & \makecell{용도지역} & \makecell{도로\newline{}교통} & \makecell{형상\newline{}지세} & \makecell{토지단가\newline{}(원/m\textsuperscript{2})} & \makecell{기준\newline{}시점} & \makecell{평가목적} \\
\hline
\makecell{a} & \makecell{S구 S동\newline{}249} & \makecell{대} & \makecell{530.0} & \makecell{상업용} & \makecell{근린상업} & \makecell{광대한면} & \makecell{부정형\newline{}평지} & \makecell{6,400,000} & \makecell{2018.\newline{}05. 01.} & \makecell{담보} \\
\hline
\makecell{b} & \makecell{S구 S동\newline{}253} & \makecell{대} & \makecell{492.0} & \makecell{상업용} & \makecell{근린상업} & \makecell{광대한면} & \makecell{사다리\newline{}평지} & \makecell{7,000,000} & \makecell{2017.\newline{}01. 01.} & \makecell{일반\newline{}거래} \\
\hline
\makecell{c} & \makecell{S구 S동\newline{}261} & \makecell{대} & \makecell{550.0} & \makecell{업무용} & \makecell{3종\newline{}일주} & \makecell{광대\newline{}소각} & \makecell{가장형\newline{}평지} & \makecell{6,500,000} & \makecell{2018.\newline{}01. 01.} & \makecell{일반\newline{}거래} \\
\hline
\end{tabular}\par}

※ 사례는 토지의 그 밖의 요인 산정시 적용하되, 사례 중 담보목적의 평가는 안전성, 환가성 등이 고려된 단가임 \par
※ 그 밖의 요인 보정방법은 표준지 기준 산정방식(비율법)을 적용함 \par
< 자료 5 > 임대사례 \par
{\small\centering
\begin{tabular}{|c|c|c|c|c|}
\hline
\makecell{소재지} & \makecell{임대면적\newline{}(m\textsuperscript{2})} & \makecell{임대보증금\newline{}(원)} & \makecell{월임대료\newline{}(원)} & \makecell{임대계약\newline{}일자} \\
\hline
\makecell{S구 S동\newline{}185-28\newline{}1층 102호} & \makecell{70} & \makecell{30,000,000} & \makecell{2,750,000\newline{}(부가세 10 \% 포함)} & \makecell{2017. 02. 02.} \\
\hline
\end{tabular}\par}

※ 보증금운용이율은 연 4 \%로 조사되어 이를 적용함 \par
※ 부가가치세(부가세)는 임차인이 환급받을 수 있음 \par
< 자료 6 > 시점수정 자료 \par
1. 지가변동률(A광역시 S구) \par
{\small\centering
\begin{tabular}{|c|c|c|c|}
\hline
\multirow{2}{*}{\makecell{기간}} & \multicolumn{2}{|c|}{\makecell{지가변동률(\%)}} & \multirow{2}{*}{\makecell{비고}} \\
\cline{2-3}
 & \makecell{상업지역} & \makecell{주거지역} &  \\
\hline
\makecell{2017. 01. 01 ∼ 2017. 12. 31} & \makecell{1.112} & \makecell{1.238} &  \\
\hline
\makecell{2018. 01. 01 ∼ 2018. 05. 31} & \makecell{1.396} & \makecell{1.574} & \makecell{2018. 05.까지 누계} \\
\hline
\makecell{2018. 05. 01 ∼ 2018. 05. 31} & \makecell{0.227} & \makecell{0.235} & \makecell{2018. 05. 당월} \\
\hline
\end{tabular}\par}

※ 2018. 06. 이후 지가변동률은 기준시점 현재 고시되지 않아 2018. 05. 지가변동률을 연장 추정하여 적용하되, 소수점 넷째자리에서 반올림하여 셋째자리까지 표시함 \par
2. 자본수익률(A광역시) \par
{\small\centering
\begin{tabular}{|c|c|c|}
\hline
\makecell{기간} & \makecell{상업용부동산\newline{}자본수익률(\%)} & \makecell{비고} \\
\hline
\makecell{2017. 01. 01 ∼ 2017. 12. 31} & \makecell{2.930} & \makecell{2017년(365일) 누계} \\
\hline
\makecell{2018. 01. 01 ∼ 2018. 03. 31} & \makecell{0.731} & \makecell{2018년 1분기} \\
\hline
\end{tabular}\par}

※ 자본수익률은 2018년도 2분기 자료가 기준시점 현재 고시되지 않아 2018년도 1분기(90일) 자본수익률을 연장 추정하여 적용하되, 소수점 넷째자리에서 반올림하여 셋째자리까지 표시함 \par
< 자료 7 > 가치형성요인 비교자료 \par
1. 접근의 편리성 \par
대상물건의 토지는 표준지 기호 1) 대비 3 \%, 표준지 기호 2) 대비 25 \%, 표준지 기호 3) 대비 5 \% 각각 우세함 \par
2. 격차율 자료 \par
1) 토지이용상황 \par
{\small\centering
\begin{tabular}{|c|c|c|c|}
\hline
\makecell{구분} & \makecell{주거용} & \makecell{상업업무용} & \makecell{주상복합용} \\
\hline
\makecell{주거용} & \makecell{1.00} & \makecell{1.30} & \makecell{1.11} \\
\hline
\makecell{상업업무용} & \makecell{0.77} & \makecell{1.00} & \makecell{0.85} \\
\hline
\makecell{주상복합용} & \makecell{0.90} & \makecell{1.17} & \makecell{1.00} \\
\hline
\end{tabular}\par}

2) 형상 \par
{\small\centering
\begin{tabular}{|c|c|c|c|c|c|c|}
\hline
\makecell{구분} & \makecell{정방형} & \makecell{가장형} & \makecell{세장형} & \makecell{사다리형} & \makecell{부정형} & \makecell{자루형} \\
\hline
\makecell{정방형} & \makecell{1.00} & \makecell{1.05} & \makecell{0.99} & \makecell{0.98} & \makecell{0.95} & \makecell{0.90} \\
\hline
\makecell{가장형} & \makecell{0.95} & \makecell{1.00} & \makecell{0.94} & \makecell{0.93} & \makecell{0.90} & \makecell{0.86} \\
\hline
\makecell{세장형} & \makecell{1.01} & \makecell{1.06} & \makecell{1.00} & \makecell{0.99} & \makecell{0.96} & \makecell{0.91} \\
\hline
\makecell{사다리형} & \makecell{1.02} & \makecell{1.07} & \makecell{1.01} & \makecell{1.00} & \makecell{0.97} & \makecell{0.92} \\
\hline
\makecell{부정형} & \makecell{1.05} & \makecell{1.11} & \makecell{1.04} & \makecell{1.03} & \makecell{1.00} & \makecell{0.95} \\
\hline
\makecell{자루형} & \makecell{1.11} & \makecell{1.17} & \makecell{1.10} & \makecell{1.09} & \makecell{1.06} & \makecell{1.00} \\
\hline
\end{tabular}\par}

3) 도시ㆍ군계획시설 \par
{\small\centering
\begin{tabular}{|c|c|c|c|}
\hline
\makecell{구분} & \makecell{일반} & \makecell{도로} & \makecell{공원} \\
\hline
\makecell{일반} & \makecell{1.00} & \makecell{0.85} & \makecell{0.60} \\
\hline
\makecell{도로} & \makecell{1.18} & \makecell{1.00} & \makecell{0.71} \\
\hline
\makecell{공원} & \makecell{1.67} & \makecell{1.42} & \makecell{1.00} \\
\hline
\end{tabular}\par}

3. 임대사례 개별요인 비교자료 \par
대상물건의 건물 중 1층 101호는 임대사례 대비 외부요인 25 \% 열세, 건물요인 10 \% 열세, 호별요인 9 \% 열세, 기타요인은 동일함 \par
4. 대상물건, 표준지, 평가사례, 임대사례 등은 모두 지역요인이 같거나 유사하고, 상기 제시자료 외의 개별요인은 모두 대등한 것으로 판단됨 \par
< 자료 8 > 재조달원가 자료 \par
1. 투자자 乙은 대상물건의 인근토지에 신축을 통한 투자를 계획하고 있음 \par
2. 건축구조 등: 철근콘크리트조 슬래브지붕 지하 1층, 지상 7층 상업용 건물 \par
3. 건물면적: 지하 1층 300 m\textsuperscript{2}, 지상 1층 180 m\textsuperscript{2}, 지상 2층∼7층 각각 250 m\textsuperscript{2} \par
4. 건축기간: 1년(건축설계 및 허가 2개월, 공사기간 10개월) \par
5. 공사비 등 투자계획 \par
1) 기준시점일(2018. 07. 01.) 현재 도급계약금액은 20억원임 \par
2) 기준시점에 5억원을 지급하고, 건물준공시점에 건물을 담보로 은행에서 잔액을 대출받아 지급함. 대출조건은 대출기간 10년, 대출이자율 연 4 \%, 매년 원리금 균등분할상환임 \par
3) 乙은 건물준공시 5년간 임대예정이고, 5년 임대기간이 종료되는 시점에 임차인에게 해당 건물을 처분하면서 대출잔금을 일시상환하려고 함 \par
4) 시장이자율은 연 5 \%임 \par
6. 대상물건 중 건물의 재조달원가는 상기조건을 고려한 건축비의 현가를 산정하여 적용하고, 내용연수는 50년, 감가수정은 정액법을 적용함 \par
7. 자본환원표 \par
1) 이자율 연 4 \% \par
{\small\centering
\begin{tabular}{|c|c|c|c|c|c|c|}
\hline
\makecell{기간\newline{}(년)} & \makecell{일시불\newline{}내가계수} & \makecell{연금\newline{}내가계수} & \makecell{감채기금\newline{}계수} & \makecell{일시불\newline{}현가계수} & \makecell{연금\newline{}현가계수} & \makecell{저당상수} \\
\hline
\makecell{1} & \makecell{1.040000} & \makecell{1.000000} & \makecell{1.000000} & \makecell{0.961538} & \makecell{0.961538} & \makecell{1.040000} \\
\hline
\makecell{5} & \makecell{1.216653} & \makecell{5.416323} & \makecell{0.184627} & \makecell{0.821927} & \makecell{4.451822} & \makecell{0.224627} \\
\hline
\makecell{10} & \makecell{1.480244} & \makecell{12.006107} & \makecell{0.083291} & \makecell{0.675564} & \makecell{8.110896} & \makecell{0.123291} \\
\hline
\end{tabular}\par}

2) 이자율 연 5 \% \par
{\small\centering
\begin{tabular}{|c|c|c|c|c|c|c|}
\hline
\makecell{기간\newline{}(년)} & \makecell{일시불\newline{}내가계수} & \makecell{연금\newline{}내가계수} & \makecell{감채기금\newline{}계수} & \makecell{일시불\newline{}현가계수} & \makecell{연금\newline{}현가계수} & \makecell{저당상수} \\
\hline
\makecell{1} & \makecell{1.050000} & \makecell{1.000000} & \makecell{1.000000} & \makecell{0.952381} & \makecell{0.952381} & \makecell{1.050000} \\
\hline
\makecell{5} & \makecell{1.276282} & \makecell{5.525631} & \makecell{0.180975} & \makecell{0.783526} & \makecell{4.329477} & \makecell{0.230975} \\
\hline
\makecell{10} & \makecell{1.628895} & \makecell{12.577893} & \makecell{0.079505} & \makecell{0.613913} & \makecell{7.721735} & \makecell{0.129505} \\
\hline
\end{tabular}\par}

< 자료 9 > 효용지수자료 \par
1. 층별효용지수 \par
{\small\centering
\begin{tabular}{|c|c|c|c|c|c|c|}
\hline
\multirow{2}{*}{\makecell{층별}} & \makecell{지하 1층} & \makecell{1층} & \makecell{2층} & \makecell{3층} & \makecell{4층} & \makecell{5층 이상} \\
\cline{2-7}
 & \makecell{45} & \makecell{100} & \makecell{52} & \makecell{46} & \makecell{44} & \makecell{42} \\
\hline
\end{tabular}\par}

※ 대상물건의 전유면적 기준 층별효용적수의 합계는 69,166임 \par
2. 호별효용지수 \par
{\small\centering
\begin{tabular}{|c|c|c|c|}
\hline
\multirow{2}{*}{\makecell{호별}} & \makecell{101호} & \makecell{102호} & \makecell{103호} \\
\cline{2-4}
 & \makecell{100} & \makecell{95} & \makecell{90} \\
\hline
\end{tabular}\par}

※ 대상물건의 전유면적 기준 1층 호별효용적수의 합계는 17,795임 \par
< 자료 10 > 기타 사항 \par
1. 대상물건은 최유효이용으로 판단됨 \par
2. 임대사례의 임대료는 신규계약된 정상임대료로 판단되며, 감정평가 대상인 1층 101호의 임대료도 연간실질임대료로 산정함 \par
3. 1층 101호의 기초가액은 층별ㆍ호별효용비율에 의한 배분방법을 적용하여 산정함 \par
4. 적산법에 적용되는 기대이율은 5 \%, 필요제경비는 순임대료(기대수익)의 7 \%임 \par
5. 요인비교치는 소수점 셋째자리에서 반올림하여 둘째자리까지 표시함 \par
6. 효용비율은 소수점 넷째자리에서 반올림하여 셋째자리까지 표시함 \par
\end{longquestion}
